\chapter{Registro de riesgos}
\label{ap:riesgos}

Este apéndice recoge el registro de riesgos que guió la ejecución del trabajo. Se incluye
por dos motivos. El primero es que varias decisiones de diseño de los capítulos anteriores
no se entienden sin él: no se eligieron por ser las mejores en abstracto, sino por ser las
que sobrevivían a un riesgo concreto. El segundo es que seis de los dieciséis riesgos
\emph{no} estaban previstos: se detectaron midiendo, no planificando, y esa es en sí misma
una observación sobre cómo se planifica un trabajo de esta naturaleza.

\section{Riesgos previstos en la planificación inicial}

Los diez de la tabla~\ref{tab:riesgos-previstos} se identificaron antes de empezar a
medir. Dos de ellos se materializaron, y que estuvieran previstos es lo que evitó que
descarrilaran el trabajo.

\begin{table}[h]
\centering
\footnotesize
\setlength{\tabcolsep}{4pt}
\begin{tabular}{|c|p{5.1cm}|p{5.6cm}|c|}
\hline
\textbf{Id} & \textbf{Riesgo} & \textbf{Mitigación aplicada} & \textbf{Estado} \\
\hline
R1 & No existe corpus educativo en español con transcripción de referencia fiable y de
     acceso abierto & Dos planes aplicados: corpus públicos con el desajuste de dominio
     medido y declarado, y subtítulos existentes verificados a mano. El tercero, la
     solicitud institucional del corpus objetivo, queda redactado y dirigido pero su
     tramitación excede el plazo & Materializado \\
\hline
R2 & Entrenar sobre tarjeta gráfica AMD puede resultar inviable & Verificar en la primera
     fase que la tarjeta \emph{entrena}, con un ciclo mínimo, antes de comprometer la
     técnica que lo necesita & Mitigado \\
\hline
R3 & Datos insuficientes para el ajuste fino & Última en el orden de ejecución y segunda en
     el de recorte; no comprometerla hasta resolver R1 y R2 & No materializado \\
\hline
R4 & Latencia insuficiente para operar en directo & Instrumentar la medida desde el primer
     prototipo y reportar la latencia real aunque no cumpla & Mitigado \\
\hline
R5 & El doble núcleo duplica el trabajo & Tratar la aplicación como vehículo de
     demostración y banco de pruebas, no como contribución con validación propia &
     Materializado \\
\hline
R6 & La comparativa no arroja diferencias concluyentes & Protocolo congelado de antemano
     para que un resultado negativo siga siendo defendible; reportar variabilidad y no solo
     medias & Materializado \\
\hline
R7 & Redacción tardía & Escribir desde la primera fase, en borradores parciales & Mitigado \\
\hline
R8 & Resultados no reproducibles & Una configuración por experimento, semillas fijadas,
     procedencia en cada salida, tablas generadas y no copiadas & Mitigado \\
\hline
R9 & Tratamiento de datos personales & No grabar material propio; el corpus empleado tiene
     licencia de investigación. La restricción de red la impone la aplicación & Mitigado \\
\hline
R10 & Disponibilidad personal & Planificación en fases relativas y puntos de recorte
      explícitos & Mitigado \\
\hline
\end{tabular}
\caption{Riesgos identificados durante la planificación.}
\label{tab:riesgos-previstos}
\end{table}

R1 se materializó en su forma más costosa: el corpus objetivo no se obtuvo, y el trabajo se
ejecutó sobre sustitutos. R6 también, y de ahí que dos de las tres técnicas evaluadas no
produzcan efecto concluyente. Que ambos estuvieran previstos es lo que permitió que su
materialización no descarrilara el trabajo: el criterio de significación y la política de
marcar como no concluyente lo que no alcanza potencia estaban fijados antes de saber qué
iban a dar los datos.

\section{Riesgos detectados midiendo}

Los seis siguientes no aparecían en la planificación. Se documentan aparte porque comparten
un rasgo, visible en la tabla~\ref{tab:riesgos-emergentes}: ninguno se manifestó como un
problema, sino como un resultado creíble.

\begin{table}[h]
\centering
\footnotesize
\setlength{\tabcolsep}{4pt}
\begin{tabular}{|c|p{5.1cm}|p{5.6cm}|c|}
\hline
\textbf{Id} & \textbf{Riesgo} & \textbf{Mitigación aplicada} & \textbf{Estado} \\
\hline
R11 & La calidad de la transcripción de referencia varía entre corpus y se confunde con
      rendimiento del modelo & Auditoría manual de una muestra antes de adoptar cualquier
      corpus; reportar el error con y sin normalización de tildes como cota superior e
      inferior & Mitigado \\
\hline
R12 & Efecto suelo: sobre voz leída el modelo sin adaptar ya acierta casi todo y no queda
      margen donde las técnicas puedan demostrar nada & Trasladar la comparativa al material
      de habla espontánea, donde el error base deja margen & Mitigado \\
\hline
R13 & Desajuste de variedad dialectal entre corpus y destino & Medido: a registro
      equivalente la diferencia entre variedades es de dos décimas, frente a los veinte
      puntos que separa el registro. Queda descartado como factor & Descartado \\
\hline
R14 & La configuración de decodificación es un factor de confusión que mueve el error hasta
      trece puntos & Congelarla e incorporarla a la identidad de cada resultado & Mitigado \\
\hline
R15 & Decisiones de alcance bloqueadas a la espera de validación externa & Tomar decisiones
      provisionales documentadas y marcadas como revisables, con cola viva de preguntas
      pendientes; avanzar mientras tanto en lo que no depende de terceros & Materializado \\
\hline
R16 & El modelo base envejece durante el desarrollo y la elección deja de estar justificada
      & Medirlo en lugar de argumentarlo: contraste con una arquitectura distinta sobre 800
      fragmentos, y catálogo de modelos auditados que se revisa igual que el de corpus &
      Mitigado \\
\hline
\end{tabular}
\caption{Riesgos detectados durante la ejecución, ausentes de la planificación inicial.}
\label{tab:riesgos-emergentes}
\end{table}

R14 y R16 merecen un comentario conjunto, porque ilustran los dos extremos del mismo
problema. R14 era invisible desde dentro: dos ejecuciones del mismo modelo sobre el mismo
corpus daban cifras incomparables y nada lo delataba. R16 era invisible desde fuera: los
resultados seguían siendo correctos, y lo que envejecía era su premisa. Ninguno de los dos
habría aparecido en una revisión del código o de los resultados; el primero exigió comparar
ejecuciones que no debían diferir, y el segundo mirar fuera de la familia de modelos
elegida.

R15 condicionó el método de trabajo más que ningún otro. La norma adoptada ante cualquier
decisión de alcance que no pudiera cerrarse en el momento fue la misma: elegir una opción
razonable por defecto, dejar constancia escrita del motivo en un registro de decisiones y
marcarla como revisable, en lugar de detener el trabajo. El coste de esa política es que
varias decisiones quedan documentadas y argumentadas pero pendientes de validación; el
beneficio es que ninguna de ellas bloqueó el avance del resto.
